\chapter{Las interacciones fundamentales}\label{ch:g11:fundamental-interactions}

Frota un globo contra el pelo y recogerá trocitos de papel ---
venciendo, con unos pocos centímetros cuadrados de goma, al tirón
gravitatorio del planeta entero. Esa victoria tan fácil es la razón de
que los \omterm{def:g11:fundamental-interactions:ladder}{átomos} se sostengan y de que los \omterm{def:g11:fundamental-interactions:ladder}{núcleos} pesados acaben por
romperse (\cref{ch:g11:nucleus-radioactivity}). El universo funciona con
exactamente cuatro interacciones; este capítulo las presenta, las pesa
unas contra otras y asigna a cada una su planta.

\section{La materia, del quark a las galaxias}

\begin{definition}[La escalera de las estructuras]\label{def:g11:fundamental-interactions:ladder}
La materia se construye por plantas, cada una ensamblada con la de
abajo: tres \emph{quarks}\index{quark} (puntuales, de menos de
\qty{e-18}{m}) se unen en un \emph{nucleón}\index{nucleón} --- protón o
neutrón, de unos \qty{e-15}{m}; los nucleones se apiñan en un
\emph{núcleo}\index{núcleo} (de \qty{e-15}{m} a \qty{e-14}{m}); un
núcleo más su nube de electrones es un \emph{átomo}\index{átomo}
(\qty{e-10}{m}); los átomos se unen en \emph{moléculas}\index{molécula}
y cristales (de \qty{e-9}{m} en adelante), que se ensamblan en objetos
cotidianos (\qty{1}{m}), planetas (\qty{e7}{m}), sistemas planetarios
(\qty{e11}{m}), galaxias (\qty{e21}{m}) y el universo observable
(\qty{e26}{m}).
\end{definition}

\begin{definition}[Las cuatro interacciones]\label{def:g11:fundamental-interactions:four}
Una \emph{interacción fundamental}\index{interacción fundamental} es una
\omterm{def:g10:inertia:force}{fuerza} que no se reduce a otras \omterm{def:g10:inertia:force}{fuerzas}; hay exactamente cuatro: la
\emph{gravitación}\index{gravitación}, la \emph{interacción
electromagnética}\index{interacción electromagnética}, la \omterm{def:g11:fundamental-interactions:strong}{interacción
fuerte} y la \omterm{def:g11:fundamental-interactions:weak}{interacción débil}.
\end{definition}

\begin{remark}\label{rem:g11:fundamental-interactions:empty}
Las plantas están separadas por vacíos: un \omterm{def:g11:fundamental-interactions:ladder}{átomo} es $10^{5}$ veces más
ancho que su \omterm{def:g11:fundamental-interactions:ladder}{núcleo} --- si se agranda el \omterm{def:g11:fundamental-interactions:ladder}{núcleo} hasta una canica de
\qty{1}{cm}, la nube de electrones mide un kilómetro. Algo tiene que
mantener unidas esas plantas tan poco pobladas, y casi nunca es la
gravedad.
\end{remark}

\begin{omfigure}
\begin{tikzpicture}[font=\small, x=0.24cm]
  \draw[-Stealth] (-1,0) -- (46.5,0);
  \foreach \x/\l/\d in {0/{10^{-18}}/2pt, 3/{10^{-15}}/13pt,
                        8/{10^{-10}}/2pt, 18/{10^{0}}/2pt,
                        25/{10^{7}}/2pt, 29/{10^{11}}/13pt,
                        39/{10^{21}}/2pt, 44/{10^{26}}/13pt}
    \draw (\x,0.1) -- (\x,-0.1) node[below=\d] {$\l$};
  \foreach \x/\n in {0/quark, 3/núcleo, 8/átomo, 18/humano, 25/Tierra,
                     29/{Sol--Tierra}, 39/galaxia, 44/universo}{
    \fill[omDef] (\x,0) circle (1.6pt);
    \node[rotate=45, anchor=west, inner sep=1.5pt] at (\x,0.24) {\n};}
\end{tikzpicture}

{\small La escalera de las estructuras en un eje de potencias de diez:
cuarenta y cuatro décadas separan el \omterm{def:g11:fundamental-interactions:ladder}{quark} del universo observable.
Tamaños en \omterm{def:g10:orders-of-magnitude:unit}{metros}.}
\end{omfigure}

\section{La carga eléctrica}

\begin{definition}[Carga eléctrica]\label{def:g11:fundamental-interactions:charge}
La \emph{carga eléctrica}\index{carga eléctrica} es la propiedad de la
materia sobre la que actúa la \omterm{def:g11:fundamental-interactions:four}{interacción electromagnética}, igual que la
masa es la propiedad sobre la que actúa la \omterm{def:g11:fundamental-interactions:four}{gravitación}. Se mide en
\emph{culombios}\index{culombio} (\unit{C}) y viene en \emph{dos
signos}, positivo y negativo, que se cancelan al sumarse; las cargas del
mismo signo se repelen y las de signos distintos se atraen.
\end{definition}

\begin{definition}[Carga elemental]\label{def:g11:fundamental-interactions:elementary}
La carga es granulada: toda carga observada es un múltiplo entero de la
\emph{carga elemental}\index{carga elemental} $e = \qty{1.60e-19}{C}$.
El protón lleva $+e$, el electrón $-e$ y el neutrón $0$; una carga $q$
contiene $N = \abs{q}/e$ cargas elementales.
\end{definition}

\begin{proposition}[Conservación de la carga]\label{prop:g11:fundamental-interactions:conservation}
La carga total de un sistema aislado no cambia nunca: la carga no se
crea ni se destruye, solo se transfiere --- en la práctica, mediante
electrones, los portadores más ligeros.
\end{proposition}
\admitted

\begin{remark}\label{rem:g11:fundamental-interactions:conservation}
Nunca se ha observado ninguna violación; la razón profunda, una simetría
del electromagnetismo, se da en los volúmenes universitarios.
\end{remark}

\begin{definition}[Electrización por frotamiento y por contacto]\label{def:g11:fundamental-interactions:charging}
Frotar dos aislantes arranca electrones de uno y los deposita en el
otro, dejando cargas opuestas de igual módulo (\emph{electrización por
frotamiento}\index{electrización por frotamiento}); tocar un objeto
neutro con uno cargado reparte la carga (\emph{electrización por
contacto}\index{electrización por contacto}). Solo se mueven los
electrones; la carga total se conserva.
\end{definition}

\begin{example}[Contar electrones]\label{ex:g11:fundamental-interactions:balloon}
Un globo frotado contra el pelo adquiere $q = \qty{-1.6e-8}{C}$: ha
\emph{ganado} $N = \num{1.6e-8} / \num{1.6e-19} = \num{1.0e11}$
electrones, y el pelo queda con exactamente $+\qty{1.6e-8}{C}$.
\end{example}

\section{La ley de Coulomb}

\begin{theorem}[Ley de Coulomb]\label{thm:g11:fundamental-interactions:coulomb}
Dos cargas puntuales $q_1$ y $q_2$ separadas una distancia $d$ ejercen
una sobre otra \omterm{def:g10:inertia:force}{fuerzas} dirigidas según la recta que las une, de módulos
iguales y sentidos opuestos --- repulsivas si los signos son iguales y
atractivas si son distintos --- con
\[
F = k\,\frac{\abs{q_1 q_2}}{d^{2}},
\qquad k = \qty{8.99e9}{N.m^2/C^2}.
\]
\end{theorem}
\admitted

\begin{remark}\label{rem:g11:fundamental-interactions:torsion}
Coulomb estableció la ley en 1785 midiendo el retorcimiento de una fina
fibra de torsión; por qué el \omterm{def:g10:orders-of-magnitude:scinot}{exponente} es exactamente $2$ se responde en
los volúmenes universitarios.
\end{remark}

\begin{omfigure}
\begin{tikzpicture}[font=\small]
  \draw[omDef, thick, fill=omDef!20] (0,0) circle (0.3) node[black] {$+$};
  \draw[omDef, thick, fill=omDef!20] (2.6,0) circle (0.3) node[black] {$+$};
  \draw[-Stealth, omThm, very thick] (-0.3,0) -- (-1.5,0)
    node[midway, above] {$\vect F_{2/1}$};
  \draw[-Stealth, omThm, very thick] (2.9,0) -- (4.1,0)
    node[midway, above] {$\vect F_{1/2}$};
\end{tikzpicture}\qquad
\begin{tikzpicture}[font=\small]
  \draw[omDef, thick, fill=omDef!20] (0,0) circle (0.3) node[black] {$+$};
  \draw[omProp, thick, fill=omProp!25] (2.6,0) circle (0.3) node[black] {$-$};
  \draw[-Stealth, omThm, very thick] (0.3,0) -- (1.1,0)
    node[midway, above] {$\vect F_{2/1}$};
  \draw[-Stealth, omThm, very thick] (2.3,0) -- (1.5,0)
    node[midway, above] {$\vect F_{1/2}$};
\end{tikzpicture}

{\small Las cargas del mismo signo se repelen (izquierda) y las de
signos distintos se atraen (derecha); las dos \omterm{def:g10:inertia:force}{fuerzas} tienen siempre
módulos iguales y sentidos opuestos.}
\end{omfigure}

\begin{remark}[Un gemelo formal de la gravitación]\label{rem:g11:fundamental-interactions:twin}
La ley de Coulomb es, símbolo a símbolo, la ley de \omterm{def:g11:fundamental-interactions:four}{gravitación} universal
(\cref{ch:g10:universal-gravitation}) con cargas en lugar de masas:
\begin{center}
\begin{tabular}{lll}
fuente & masa $m$ & carga $q$ \\
ley & $F = G\,\dfrac{m_1 m_2}{d^2}$ & $F = k\,\dfrac{\abs{q_1 q_2}}{d^2}$ \\[1.5ex]
constante & $G = \qty{6.67e-11}{N.m^2/kg^2}$ & $k = \qty{8.99e9}{N.m^2/C^2}$ \\
signo & siempre atractiva & atractiva o repulsiva \\
\end{tabular}
\end{center}
Las dos diferencias gobiernan todo lo que sigue: la electricidad es
abrumadoramente más fuerte, y solo ella puede cancelarse a sí misma.
\end{remark}

\begin{example}[Dos protones, las dos fuerzas a la vez]\label{ex:g11:fundamental-interactions:protons}
Dos protones ($m_p = \qty{1.67e-27}{kg}$) están a
$d = \qty{1.0e-15}{m}$ --- vecinos dentro de un \omterm{def:g11:fundamental-interactions:ladder}{núcleo}. Repulsión
eléctrica:
\[
F_E = \num{8.99e9} \times
\frac{(\num{1.60e-19})^{2}}{(\num{1.0e-15})^{2}}
\approx \qty{2.3e2}{N}
\]
--- el \omterm{def:g10:universal-gravitation:weight}{peso} de una maleta de $23$ kg, soportado por una partícula de
$10^{-27}$ \omterm{def:g10:orders-of-magnitude:unit}{kilogramos}. Atracción gravitatoria:
$F_G = \num{6.67e-11} \times (\num{1.67e-27})^{2} /
(\num{1.0e-15})^{2} \approx \qty{1.9e-34}{N}$. El cociente
$F_E / F_G = k e^{2} / (G m_p^{2}) \approx \num{1.2e36}$ no depende de
la distancia ($d^{2}$ se cancela): entre partículas elementales, la
gravedad es $10^{36}$ veces demasiado débil para contar --- a
\emph{cualquier} distancia.
\end{example}

\section{Las interacciones fuerte y débil}

El ejemplo deja un escándalo: un \omterm{def:g11:fundamental-interactions:ladder}{núcleo} de helio sostiene dos protones
que se repelen con decenas de \omterm{def:g10:inertia:force}{newtons} --- y sin embargo el helio existe.

\begin{definition}[La interacción fuerte]\label{def:g11:fundamental-interactions:strong}
La \emph{interacción fuerte}\index{interacción fuerte} liga los \omterm{def:g11:fundamental-interactions:ladder}{quarks}
en nucleones y los nucleones en \omterm{def:g11:fundamental-interactions:ladder}{núcleos}. A \qty{e-15}{m} es atractiva y
mucho más intensa que la repulsión eléctrica (miles de \omterm{def:g10:inertia:force}{newtons} entre
nucleones); más allá de unos pocos $10^{-15}$ \omterm{def:g10:orders-of-magnitude:unit}{metros} se desvanece: su
\emph{alcance}\index{alcance (de una interacción)} es de unos
\qty{e-15}{m}. Agarra por igual a protones y neutrones e ignora la
carga.
\end{definition}

\begin{omfigure}
\begin{tikzpicture}[font=\small]
  \draw[omDef, thick, fill=omDef!20] (0,0) circle (0.34) node[black] {p};
  \draw[omDef, thick, fill=omDef!20] (3.2,0) circle (0.34) node[black] {p};
  \draw[-Stealth, omThm, thick] (-0.34,0) -- (-1.7,0)
    node[midway, above] {$F_E \approx \qty{230}{N}$};
  \draw[-Stealth, omThm, thick] (3.54,0) -- (4.9,0) node[midway, above] {$F_E$};
  \draw[-Stealth, omMeth, very thick] (0.2,0.55) -- (1.45,0.55)
    node[midway, above] {fuerte};
  \draw[-Stealth, omMeth, very thick] (3.0,0.55) -- (1.75,0.55)
    node[midway, above] {fuerte};
  \draw[Stealth-Stealth] (0,-0.65) -- (3.2,-0.65)
    node[midway, below] {$d = \qty{1.0e-15}{m}$};
\end{tikzpicture}

{\small El tira y afloja dentro de un \omterm{def:g11:fundamental-interactions:ladder}{núcleo}: a $10^{-15}$~m la
atracción fuerte (miles de \omterm{def:g10:inertia:force}{newtons}) le gana a la repulsión eléctrica.}
\end{omfigure}

\begin{remark}[Por qué existen los núcleos --- y por qué los pesados ceden]\label{rem:g11:fundamental-interactions:nuclei}
Las dos \omterm{def:g10:inertia:force}{fuerzas} del \omterm{def:g11:fundamental-interactions:ladder}{núcleo} escalan de forma distinta. El pegamento
fuerte es de corto \omterm{def:g11:fundamental-interactions:strong}{alcance}: un \omterm{def:g11:fundamental-interactions:ladder}{nucleón} solo se liga a sus pocos vecinos
situados dentro de \qty{e-15}{m}. La repulsión eléctrica es de largo
\omterm{def:g11:fundamental-interactions:strong}{alcance}: cada protón empuja a \emph{todos los demás} protones. A medida
que los \omterm{def:g11:fundamental-interactions:ladder}{núcleos} crecen, la repulsión se acumula más deprisa que el
pegamento: los \omterm{def:g11:fundamental-interactions:ladder}{núcleos} pesados necesitan neutrones de más (pegamento sin
carga) y, más allá de unos $80$ protones, ni siquiera eso basta --- los
\omterm{def:g11:fundamental-interactions:ladder}{núcleos} más pesados viven de prestado
(\cref{ch:g11:nucleus-radioactivity}).
\end{remark}

\begin{definition}[La interacción débil]\label{def:g11:fundamental-interactions:weak}
La \emph{interacción débil}\index{interacción débil}, de \omterm{def:g11:fundamental-interactions:strong}{alcance} aún más
corto, permite que un neutrón se convierta en protón (y al revés): es
ella la que impulsa la radiactividad $\beta$ del
\cref{ch:g11:nucleus-radioactivity}.
\end{definition}

\section{Quién manda en cada planta}

\begin{method}[Auditar una escala]\label{met:g11:fundamental-interactions:audit}
Para saber qué interacción manda en una estructura de tamaño $d$, hay
que mirar el \omterm{def:g11:fundamental-interactions:strong}{alcance} y la cancelación:
\begin{enumerate}
  \item $d \leq \qty{e-14}{m}$: manda la \emph{\omterm{def:g11:fundamental-interactions:strong}{interacción fuerte}} ---
        le gana a la repulsión eléctrica, y la gravedad queda $10^{36}$
        veces fuera.
  \item del \omterm{def:g11:fundamental-interactions:ladder}{átomo} a las montañas (de $\qty{e-10}{m}$ a $\qty{e4}{m}$):
        la \omterm{def:g10:inertia:force}{fuerza} fuerte queda fuera de \omterm{def:g11:fundamental-interactions:strong}{alcance}; manda la
        \emph{\omterm{def:g11:fundamental-interactions:four}{interacción electromagnética}} --- enlaces, cohesión,
        \omterm{def:g10:inertia:inventory}{rozamiento}, contacto.
  \item planetas y más allá ($d \geq \qty{e6}{m}$): la materia es neutra
        con una precisión fantástica, así que la electricidad se cancela
        a sí misma; la masa tiene un solo signo y siempre se suma ---
        manda la \emph{\omterm{def:g11:fundamental-interactions:four}{gravitación}} en el cielo.
\end{enumerate}
La \omterm{def:g11:fundamental-interactions:weak}{interacción débil} no ocupa ninguna planta: no liga nada, pero arbitra
las transformaciones (la desintegración $\beta$).
\end{method}

\begin{omfigure}
\begin{tikzpicture}[font=\small, x=0.26cm]
  \draw[-Stealth] (-1,0) -- (41.5,0);
  \foreach \x/\l in {3/{10^{-15}}, 8/{10^{-10}}, 18/{10^{0}},
                     24/{10^{6}}, 39/{10^{21}}}
    \draw (\x,0.1) -- (\x,-0.1) node[below] {$\l$};
  \draw[omProp, fill=omProp!35] (3,0.3) rectangle (4,0.75);
  \node[anchor=west, omProp] at (4.4,0.52) {fuerte};
  \draw[omDef, fill=omDef!25] (8,0.95) rectangle (22,1.4);
  \node[omDef] at (15,1.17) {electromagnética};
  \draw[omThm, fill=omThm!25] (24,1.6) rectangle (39,2.05);
  \node[omThm] at (31.5,1.82) {gravitación};
\end{tikzpicture}

{\small Quién manda dónde: la fuerte en los \omterm{def:g11:fundamental-interactions:ladder}{núcleos}, la
electromagnética del \omterm{def:g11:fundamental-interactions:ladder}{átomo} a las montañas y la \omterm{def:g11:fundamental-interactions:four}{gravitación} en los
planetas y más allá. Tamaños en \omterm{def:g10:orders-of-magnitude:unit}{metros}.}
\end{omfigure}

\section{Ejercicios}

\begin{exercise}[$\star$]\label{exo:g11:fundamental-interactions:1}
Da el \omterm{def:g10:orders-of-magnitude:oom}{orden de magnitud} de cada tamaño y su planta en la escalera:
protón \qty{1.7e-15}{m}; \omterm{def:g11:fundamental-interactions:ladder}{molécula} de agua \qty{2.8e-10}{m}; grano de
arena \qty{2e-4}{m}; Tierra \qty{1.3e7}{m}; Vía Láctea
\qty{9.5e20}{m}.
\end{exercise}

\begin{exercise}[$\star$]\label{exo:g11:fundamental-interactions:2}
Un globo frotado contra el pelo lleva $q = \qty{-4.8e-9}{C}$. ¿Ha ganado
o ha perdido electrones, y cuántos? ¿Qué carga queda después en el pelo,
y qué ley lo dice?
\end{exercise}

\begin{exercise}[$\star$]\label{exo:g11:fundamental-interactions:3}
Las cargas $q_1 = \qty{+2.0e-6}{C}$ y $q_2 = \qty{-3.0e-6}{C}$ están a
\qty{30}{cm}. Calcula la \omterm{def:g10:inertia:force}{fuerza}; ¿atractiva o repulsiva? Compara las
\omterm{def:g10:inertia:force}{fuerzas} sobre $q_1$ y sobre $q_2$.
\end{exercise}

\begin{exercise}[$\star$]\label{exo:g11:fundamental-interactions:4}
Dos cargas se atraen con una \omterm{def:g10:inertia:force}{fuerza} $F_0$. ¿En qué se convierte si se
duplica la distancia; si la distancia se divide entre $3$; si se duplica
una carga; si se duplican las dos cargas \emph{y} la distancia?
\end{exercise}

\begin{exercise}[$\star$]\label{exo:g11:fundamental-interactions:5}
Nombra la interacción responsable de: (a) la Luna orbitando la Tierra;
(b) la cohesión de un cristal de sal; (c) la cohesión de un \omterm{def:g11:fundamental-interactions:ladder}{núcleo} de
helio; (d) la desintegración $\beta$ del carbono-14; (e) el globo del
\cref{exo:g11:fundamental-interactions:2} pegado a una pared.
\end{exercise}

\begin{exercise}[$\star\star$]\label{exo:g11:fundamental-interactions:6}
Los dos protones de un \omterm{def:g11:fundamental-interactions:ladder}{núcleo} de helio están a $d = \qty{2.0e-15}{m}$
($m_p = \qty{1.67e-27}{kg}$). Calcula su repulsión eléctrica, su
atracción gravitatoria y el cociente. ¿Qué mantiene entero el \omterm{def:g11:fundamental-interactions:ladder}{núcleo}?
\end{exercise}

\begin{exercise}[$\star\star$]\label{exo:g11:fundamental-interactions:7}
En un \omterm{def:g11:fundamental-interactions:ladder}{átomo} de hidrógeno, el electrón ($m_e = \qty{9.11e-31}{kg}$) está
a $d = \qty{5.3e-11}{m}$ del protón ($m_p = \qty{1.67e-27}{kg}$).
Calcula las \omterm{def:g10:inertia:force}{fuerzas} eléctrica y gravitatoria entre ellos y su cociente.
¿Qué mantiene unidos a los \omterm{def:g11:fundamental-interactions:ladder}{átomos}?
\end{exercise}

\begin{exercise}[$\star\star$]\label{exo:g11:fundamental-interactions:8}
La esfera A lleva $q_A = \qty{+8.0e-9}{C}$; una esfera B idéntica está
neutra. Se ponen en contacto y después se separan \qty{10}{cm}. ¿Qué
carga lleva cada una (comprueba la conservación de la carga)? Calcula la
\omterm{def:g10:inertia:force}{fuerza} entre ellas --- ¿atractiva o repulsiva?
\end{exercise}

\begin{exercise}[$\star\star$]\label{exo:g11:fundamental-interactions:9}
Dos cargas iguales $q = \qty{1.0e-6}{C}$ se repelen con \qty{0.90}{N}:
¿a qué distancia están? ¿Dónde baja la \omterm{def:g10:inertia:force}{fuerza} a \qty{0.10}{N}?
\end{exercise}

\begin{exercise}[$\star\star$]\label{exo:g11:fundamental-interactions:10}
Dos protones de \omterm{def:g11:fundamental-interactions:ladder}{moléculas} vecinas están a $d = \qty{1.0e-10}{m}$.
Calcula su repulsión eléctrica. ¿Qué aporta la \omterm{def:g11:fundamental-interactions:strong}{interacción fuerte} a esa
distancia? ¿Qué interacción liga las \omterm{def:g11:fundamental-interactions:ladder}{moléculas}?
\end{exercise}

\begin{exercise}[$\star\star$]\label{exo:g11:fundamental-interactions:11}
Cada una de dos monedas contiene unos $\num{e23}$ electrones. Pasa un
electrón de cada \emph{millón} de una a la otra y sepáralas
\qty{1.0}{m}. Calcula la carga de cada moneda y después la \omterm{def:g10:inertia:force}{fuerza}. ¿Al
\omterm{def:g10:universal-gravitation:weight}{peso} de qué masa equivale? Concluye sobre la neutralidad de la materia
cotidiana.
\end{exercise}

\begin{exercise}[$\star\star\star$]\label{exo:g11:fundamental-interactions:12}
Un \omterm{def:g11:fundamental-interactions:ladder}{núcleo} de uranio tiene un diámetro de \qty{1.4e-14}{m} y $92$
protones. Calcula la repulsión entre dos protones situados en extremos
opuestos. ¿Puede la \omterm{def:g11:fundamental-interactions:strong}{interacción fuerte} ligar esa pareja? ¿Por qué no?
¿Cuántas parejas de protones que se repelen hay? ¿Por qué los \omterm{def:g11:fundamental-interactions:ladder}{núcleos}
pesados necesitan neutrones de más, y por qué a partir de cierto tamaño
ni siquiera los neutrones pueden salvarlos?
\end{exercise}

\begin{exercise}[$\star\star\star$]\label{exo:g11:fundamental-interactions:13}
Dos bolitas de masa \qty{1.0}{g} cuelgan de hilos de \qty{50}{cm}
atados al mismo punto. Al darles la misma carga $q$, se separan
\qty{6.0}{cm}. Cada bolita está en equilibrio bajo su \omterm{def:g10:universal-gravitation:weight}{peso}, la tensión
del hilo y la \omterm{def:g10:inertia:force}{fuerza} eléctrica: demuestra que $F = mg\tan\theta$
($\theta$: ángulo del hilo con la vertical) y que aquí
$\tan\theta \approx 0.060$. Calcula $F$, después $q$ y después el número
de \omterm{def:g11:fundamental-interactions:elementary}{cargas elementales} desplazadas.
\end{exercise}

\begin{exercise}[$\star\star\star$]\label{exo:g11:fundamental-interactions:14}
Datos: $M_E = \qty{5.97e24}{kg}$, $M_M = \qty{7.35e22}{kg}$, distancia
Tierra--Luna $d = \qty{3.84e8}{m}$. Calcula la \omterm{def:g10:universal-gravitation:force}{fuerza gravitatoria}
Tierra--Luna. Se apaga la gravedad: ¿qué cargas $+Q$ en la Tierra y $-Q$
en la Luna mantendrían la misma atracción? ¿A cuántos electrones
equivale $Q$, y cuánto pesan ($m_e = \qty{9.11e-31}{kg}$)? Coméntalo.
\end{exercise}

\begin{exercise}[$\star\star\star$]\label{exo:g11:fundamental-interactions:15}
En un cristal de sal, los iones Na$^{+}$ y Cl$^{-}$ (cargas $\pm e$)
están a $d = \qty{2.8e-10}{m}$; un ion Cl$^{-}$ tiene una masa de
\qty{5.9e-26}{kg}. Calcula la \omterm{def:g10:inertia:force}{fuerza} entre iones vecinos y el \omterm{def:g10:universal-gravitation:weight}{peso} de un
ion Cl$^{-}$; da el cociente. Explica después cómo gana la gravedad a
gran escala --- ninguna montaña supera los \qty{e4}{m} --- aunque cada
enlace le gane por quince órdenes de magnitud.
\end{exercise}

\section{Problema: la auditoría de las cuatro fuerzas}

\begin{problem}[{Problema de fin de semana --- por qué el mundo ni se
derrumba ni se desintegra: las cuatro interacciones auditadas planta por
planta, hasta la neutralidad de la materia con dos partes en
$10^{18}$}]
\label{pb:g11:fundamental-interactions:1}
Los \omterm{def:g11:fundamental-interactions:ladder}{átomos} no implosionan, los \omterm{def:g11:fundamental-interactions:ladder}{núcleos} casi siempre aguantan y la Luna
ni se estrella ni se escapa; este problema audita a quién corresponde el
mérito. Datos: $e = \qty{1.60e-19}{C}$,
$k = \qty{8.99e9}{N.m^2/C^2}$, $G = \qty{6.67e-11}{N.m^2/kg^2}$,
$m_p = \qty{1.67e-27}{kg}$, $m_e = \qty{9.11e-31}{kg}$,
$M_E = \qty{5.97e24}{kg}$, $M_M = \qty{7.35e22}{kg}$, distancia
Tierra--Luna \qty{3.84e8}{m}.

\medskip
\noindent\textbf{Parte I --- La escalera.}
\begin{enumerate}
  \item Enumera las plantas de la escalera del \omterm{def:g11:fundamental-interactions:ladder}{quark} a la galaxia, con
        un \omterm{def:g10:orders-of-magnitude:oom}{orden de magnitud} de tamaño cada una.
  \item Maqueta a escala: el \omterm{def:g11:fundamental-interactions:ladder}{núcleo} se convierte en una canica de
        \qty{1}{cm}; ¿cuánto mide entonces el \omterm{def:g11:fundamental-interactions:ladder}{átomo}?
  \item ¿Qué fracción del volumen de un \omterm{def:g11:fundamental-interactions:ladder}{átomo} ocupa su \omterm{def:g11:fundamental-interactions:ladder}{núcleo}?
  \item Un cabello tiene \qty{70}{\micro m} de grosor. ¿Cuántos \omterm{def:g11:fundamental-interactions:ladder}{átomos}
        lo abarcan?
  \item ¿Cuántas potencias de diez hay del techo del \omterm{def:g11:fundamental-interactions:ladder}{quark}
        (\qty{e-18}{m}) a una galaxia (\qty{e21}{m})?
\end{enumerate}

\medskip
\noindent\textbf{Parte II --- La planta eléctrica.}
\begin{enumerate}[resume]
  \item En un \omterm{def:g11:fundamental-interactions:ladder}{átomo} de hidrógeno, el electrón está a
        $d = \qty{5.3e-11}{m}$ del protón: calcula la \omterm{def:g10:inertia:force}{fuerza} eléctrica
        entre ellos.
  \item Calcula la \omterm{def:g10:universal-gravitation:force}{fuerza gravitatoria} entre ellos y el cociente de las
        dos. ¿Cuál mantiene unido el \omterm{def:g11:fundamental-interactions:ladder}{átomo}?
  \item ¿Qué cambia si se apaga la gravedad dentro de los \omterm{def:g11:fundamental-interactions:ladder}{átomos}? ¿Y si
        se apaga la electricidad?
  \item ¿Por qué las \omterm{def:g10:inertia:force}{fuerzas} de ``contacto'' de la vida diaria --- el
        suelo que empuja tus pies, el \omterm{def:g10:inertia:inventory}{rozamiento} --- son
        electromagnéticas disfrazadas?
  \item Cada una de dos monedas contiene unos $\num{e23}$ electrones.
        Pasa un electrón de cada \emph{mil millones} de una a la otra y
        sepáralas \qty{1.0}{m}: calcula las cargas y la \omterm{def:g10:inertia:force}{fuerza}.
        ¿Conclusión?
  \item ¿Qué interacción manda en las plantas de \qty{e-10}{m} a
        \qty{e4}{m}, y por qué no compite la \omterm{def:g11:fundamental-interactions:strong}{interacción fuerte}?
\end{enumerate}

\medskip
\noindent\textbf{Parte III --- La planta nuclear.}
\begin{enumerate}[resume]
  \item Dos protones están a \qty{1.0e-15}{m} dentro de un \omterm{def:g11:fundamental-interactions:ladder}{núcleo}:
        calcula su repulsión eléctrica.
  \item Sin nada que lo frene, ¿qué aceleración le daría esa \omterm{def:g10:inertia:force}{fuerza} a un
        protón? Compárala con $g$.
  \item Enumera las tres propiedades que necesita la \omterm{def:g11:fundamental-interactions:strong}{interacción fuerte}
        para explicar por qué existen los \omterm{def:g11:fundamental-interactions:ladder}{núcleos}, por qué también
        agarra a los neutrones y por qué los \omterm{def:g11:fundamental-interactions:ladder}{núcleos} son diminutos.
  \item En un \omterm{def:g11:fundamental-interactions:ladder}{núcleo} de uranio (diámetro \qty{1.4e-14}{m}, $92$
        protones), calcula la repulsión entre dos protones situados en
        extremos opuestos. ¿Puede la \omterm{def:g11:fundamental-interactions:strong}{interacción fuerte} ligar
        directamente esa pareja? Explica por qué la repulsión gana a
        medida que crecen los \omterm{def:g11:fundamental-interactions:ladder}{núcleos} y para qué sirven los neutrones.
  \item En una frase, como se prometió: ¿qué hace la \omterm{def:g11:fundamental-interactions:weak}{interacción débil}?
\end{enumerate}

\medskip
\noindent\textbf{Parte IV --- La planta astronómica y el veredicto.}
\begin{enumerate}[resume]
  \item Calcula la \omterm{def:g10:universal-gravitation:force}{fuerza gravitatoria} entre la Tierra y la Luna.
  \item La Tierra contiene unos $\num{1.8e51}$ protones (y otros tantos
        electrones), y la Luna unos $\num{2.2e49}$. ¿Qué fracción $f$ de
        la carga protónica de cada cuerpo, dejada sin cancelar, haría
        que la \omterm{def:g10:inertia:force}{fuerza} eléctrica igualase a la gravitatoria?
  \item ¿Por qué manda en la astronomía la \omterm{def:g11:fundamental-interactions:four}{gravitación}, la endeble de la
        pregunta 7? Dos razones.
  \item Veredicto --- una frase por planta (\omterm{def:g11:fundamental-interactions:ladder}{núcleo}, del \omterm{def:g11:fundamental-interactions:ladder}{átomo} a la
        montaña, planeta y más allá): ¿qué la mantiene unida? Responde
        después al título: ¿por qué el mundo ni se derrumba ni se
        desintegra?
\end{enumerate}
\end{problem}
